\documentclass[11pt,a4paper]{article}
\usepackage{fontspec,amsmath,amssymb,amsthm,microtype}
\setmainfont{lmroman10-regular.otf}[BoldFont=lmroman10-bold.otf,ItalicFont=lmroman10-italic.otf,BoldItalicFont=lmroman10-bolditalic.otf]
\setmonofont{lmmono10-regular.otf}[Scale=MatchLowercase]
\usepackage[margin=25mm]{geometry}
\usepackage[unicode,hidelinks]{hyperref}
\setlength{\emergencystretch}{3em}
\setlength{\parskip}{0.4em}
\setlength{\parindent}{0pt}
\allowdisplaybreaks[2]
\raggedbottom
\title{\texorpdfstring{From the actual integral equation to the exponential mild equation}{From the actual integral equation to the exponential mild equation}}
\author{LangevinLean mathematical exposition}
\date{}
\begin{document}\maketitle
\noindent Source module: \texttt{Kinetic\allowbreak{}Langevin/\allowbreak{}Mild\allowbreak{}Equation.\allowbreak{}lean}. The accompanying HTML pairs this mathematical exposition with the exact Lean statements and records the scope of the formalization.


The KLMC update freezes the force inside an exponentially weighted integral. To compare it with an exact path, that weighted formula must follow from the actual integral equation driven by the same Brownian input. The theorem below proves this bridge for any complete real normed space, including the L\ensuremath{{}^{2}} space of random Euclidean vectors. No time derivative of Brownian motion is used.


\section{\texorpdfstring{Subtract the filtered noise and apply an integrating factor}{Subtract the filtered noise and apply an integrating factor}}


Assume continuous velocity \(V\), continuous force path \(F\), and a continuous filtered input \(B\), satisfying \(V_t-V_0=-\int_0^t(\gamma V_s+F_s)ds+W_t\) and \(B_t=W_t-\gamma\int_0^tB_sds\). The conclusion is the exact identity \(V_T=e^{-\gamma T}V_0-\int_0^Te^{-\gamma(T-s)}F_sds+B_T\). In the proof, evaluation at zero first gives \(B_0=0\). Subtracting the two equations shows \(D=V-B\) solves \(D_t=V_0-\int_0^t(\gamma D_s+F_s)ds\). Continuous integrands allow the fundamental theorem of calculus on the interval interior. Lean derives \(D\!\prime=-\gamma D-F\), differentiates \(e^{\gamma t}D_t\), and obtains derivative \(-e^{\gamma t}F_t\). A second application of the fundamental theorem integrates this identity over the closed interval; the proof handles the endpoints by continuity. Finally it multiplies by \(e^{-\gamma T}\), commutes scalar multiplication with the Bochner integral, and simplifies the exponent. The hypotheses are the genuine integral equations, and the resulting mild equation is proved rather than postulated. The argument works for arbitrary real friction and includes \(T=0\).



{\raggedright\noindent Lean declarations: \texttt{damped\_\allowbreak{}integral\_\allowbreak{}eq\_\allowbreak{}mild}.\par}

\end{document}
